\begin{abcliste}
\abc
Auf die Kiste wirken die Gewichtskraft $\vec F_\mathrm{G}$ nach unten und die Federkraft $\vec F_\mathrm{F}$ nach oben. Die Feder ist gestaucht und drückt zurück zu ihrer entspannten Länge, also \result{\text{nach oben}}. Da die Kiste ruht, sind die beiden Kräfte gleich gross.
\begin{center}
\begin{tikzpicture}[>=Stealth]
 \draw[fill=brown!30] (-0.6,0) rectangle (0.6,0.8);
 \draw[->, very thick] (-0.15,0.4) -- (-0.15,-1.1) node[left] {$\vec F_\mathrm{G}$};
 \fill (-0.15,0.4) circle (0.05);
 \draw[->, very thick] (0.15,0) -- (0.15,1.5) node[right] {$\vec F_\mathrm{F}$};
 \fill (0.15,0) circle (0.05);
\end{tikzpicture}
\end{center}

\abc
Gleichgewicht: $F_\mathrm{F} = k\,\Delta x = m\,g$, also
\begin{align}
 k &= \kFF \\
   &= \frac{\mK\cdot\g}{\dx} \\
   &= \kF \approx \result{\kFP{0}{3}}
\end{align}

\abc
Wieder ruht die Kiste, und wieder muss die Feder die Gewichtskraft $m\,g = \FGP{0}{3}$ tragen. Die Feder wird also gleich stark verformt, diesmal gedehnt:
\begin{align}
 \Delta x' &= \dxsF = \dxs \approx \result{\dxsP{-2}{3}}
\end{align}
Die gedehnte Feder zieht zurück zu ihrer entspannten Länge, also zieht sie die Kiste \result{\text{nach oben}}. In beiden Fällen zeigt die Federkraft auf die Kiste nach oben; einmal drückt die gestauchte Feder von unten, einmal zieht die gedehnte Feder von oben.
\end{abcliste}
